% =====================================================================
%  Chapter 9 -- Future scope
% =====================================================================

\section{Noise as a control parameter}\label{sec:future-noise}

The growth budget predicts that the critical rate for tipping depends on the noise floor only through a
logarithm (\cref{eq:lamc}),
\begin{equation}
  \lamc \propto \frac{1}{\ln(\Asat/\delta)} .
  \label{eq:future-noise}
\end{equation}
Injecting a controlled noise voltage, for example through the input resistor, would
vary $\delta$ over several decades and test \cref{eq:future-noise} directly: a noisier circuit should
tolerate faster sweeps.

\section{Coupled oscillators and networks}\label{sec:future-networks}

Wien-bridge oscillators can be coupled resistively into networks of controllable structure, and such networks have been used to demonstrate community
switching en route to global synchronisation, on the bench and in simulation~\cite{Yadav2025}. Building
on that coupling scheme, three questions arise. In two coupled oscillators, does each critical rate shift,
and can a re-igniting oscillator rescue a neighbour that would tip? In a network, does tipping spread, and
does community structure decide where it starts and stops? In a multilayer network with one swept layer,
how much rate dependence passes to the others?
